% !TeX root = ../proyecto.tex

\begin{UseCase}{CU-04}{Identificación por enlace externo}{
	Permite al analista someter una URL periodística arbitraria que no fue capturada por el sistema automático, extraer su contenido, validar su relevancia e investigarla profundamente reutilizando el pipeline del módulo DeepInvestigator.
}
	\UCitem{Versión}{\color{Gray}
		1.0
	}\UCitem{Autor}{\color{Gray}
		Héctor Alberto Morales Martínez
	}\UCitem{Supervisa}{\color{Gray}
		Director de Trabajo Terminal
	}\UCitem{Actor}{
		Analista de datos / Trabajador social
	}\UCitem{Propósito}{\begin{Titemize}
		\Titem Incorporar al sistema casos relevantes descubiertos manualmente por el usuario navegando en la web, mitigando los falsos negativos del recolector automatizado.
	\end{Titemize}
	}\UCitem{Entradas}{\begin{Titemize}
		\Titem URL externa del artículo periodístico.
	\end{Titemize}
	}\UCitem{Origen}{\begin{Titemize}
		\Titem Formulario de ''Identificar caso'' en el Dashboard.
	\end{Titemize}
	}\UCitem{Salidas}{\begin{Titemize}
		\Titem Ficha estructurada en formato JSON (Ubicación, Víctimas, NNA Afectados, etc.) del nuevo caso.
	\end{Titemize}
	}\UCitem{Destino}{\begin{Titemize}
		\Titem Interfaz del dashboard (Identificación de caso).
		\Titem Base de datos.
	\end{Titemize}
	}\UCitem{Precondiciones}{\begin{Titemize}
		\Titem La URL debe ser accesible de forma pública.
		\Titem El servicio de Gemini Flash debe estar disponible.
	\end{Titemize}
	}\UCitem{Postcondiciones}{\begin{Titemize}
		\Titem Se incorpora una nueva noticia validada e investigada a la base de datos del sistema.
	\end{Titemize}
	}\UCitem{Errores}{\begin{Titemize}
		\Titem {\bf \hypertarget{CU-04.E1}{E1}}: La URL apunta a contenido no válido (cookies, Error 403, etc) \Trayref{CU-04}{A}.
	\end{Titemize}
	}\UCitem{Tipo}{
		Caso de uso primario
	}\UCitem{Observaciones}{
		Este módulo actúa como una vía de inyección manual que aprovecha la infraestructura de extracción y análisis ya existente en el sistema.
	}
\end{UseCase}

%--------------------------------------
\begin{UCtrayectoria}
	\UCpaso[\UCActor] Ingresa al apartado de ''Identificar caso'' en el dashboard.
	\UCpaso[\UCActor] Pega la URL de la noticia en el campo de texto y hace clic en \IUbutton{Analizar caso}.
	\UCpaso[\UCsist] Accede a la URL proporcionada mediante \texttt{StealthSession} y extrae el texto del artículo aplicando rutinas de limpieza.
	\UCpaso[\UCsist] Envía el texto extraído a Gemini Flash solicitando la generación de un título corto o la validación del contenido \ErrorRef{CU-04}{E1}{Contenido Inválido}.
	\UCpaso[\UCsist] Invoca al módulo interno \texttt{DeepInvestigator} utilizando el texto validado y el nuevo título.
	\UCpaso[\UCsist] Genera una consulta específica, busca en DuckDuckGo, hace scraping multi-fuente y sintetiza la ficha (mismos pasos que CU-03).
	\UCpaso[\UCsist] Almacena la noticia y su ficha de investigación en la base de datos.
	\UCpaso[\UCsist] Muestra los resultados estructurados en la interfaz del usuario.
\end{UCtrayectoria}

%--------------------------------------
\begin{UCtrayectoriaA}{CU-04}{A}{La extracción retorna contenido que no corresponde a una noticia válida}
	\UCpaso[\UCsist] El modelo Gemini evalúa el texto y retorna la bandera \texttt{NOTICIA\_INVALIDA}.
	\UCpaso[\UCsist] Detiene el flujo de ejecución antes de iniciar la investigación profunda.
	\UCpaso[\UCsist] Retorna un mensaje de error a la interfaz web indicando que la URL parece estar bloqueada, es un aviso de privacidad o no contiene cuerpo periodístico.
	\UCpaso[\UCActor] Lee el error y puede intentar con una URL alternativa.
	\UCpaso[] El Caso de Uso termina sin registrar información en la base de datos para no corromper los registros.
\end{UCtrayectoriaA}
